% BOOK-FILE-ID: f-3e1b30729f34
% section: 応用パターン

\section{応用パターン}\label{節：応用パターン}
  数学の問題を解く中での大原則の1つに,既知の形を作るというものがある.
  それに倣って,与えられた漸化式をいかにして\sectionref{節：基本パターン}で扱った漸化式に変形するかを考える.

  その最初として,特性方程式型を扱う.
  ここは本書で扱う漸化式の中で特段難しいものではないが,漸化式から一般項を出すという目的の中で極めて重要な発想及びその根底に存在する数学的事実とそこに至るためのいくつもの手法を扱う.
  
  \BookImport{f-e539f54ffe42}
  \BookImport{f-0140c19b880e}
  \BookImport{f-8f1004cf85dc}
  \BookImport{f-4663a7fcd21b}
  \BookImport{f-ef6f4e798ac0}
  \BookImport{f-d6b88598ed80}
  \BookImport{f-c8997aff583f}

ここまでの解法を使い分ける練習には,\sectionref*{節：解法を選ぶ}と\sectionref*{節：標準},\sectionref*{節：応用}へ進もう.
発展的な内容を先にすべて読む必要はない.\booknavref{節：演習の進め方}{「演習の進め方」}では,各パターンに対応する代表題を選べる.
